\documentclass[12pt]{article}
\usepackage{amsmath, amsfonts, amssymb}
\title{2nd semester problem sets}
\begin{document}
\maketitle
\section{Problem set 1}
\textbf{Q1.} Determine which of the following are subrings of $\mathbb{Q}$
\begin{enumerate}
    \item[a.] rational numbers with odd denominator
    \item[b.] rational numbers with even denominator
    \item[c.] nonnegative rationals
    \item[d.] squares of rationals
    \item[e.] rationals with odd numerator
    \item[f.] rationals with even numerator    
\end{enumerate}
\textbf{Soln: } Let's check option (a). Let $\mathbb{Q'}=\biggl\{\dfrac{a}{b}\mid a \in \mathbb{Z}, b=2k+1, \text{for some integer } k\biggr\}$ Suppose $\dfrac{a}{b}$ and $\dfrac{c}{d}$ are two elements of the set $\mathbb{Q'}$. Then we have that $b=2k+1$ and $d=2m+1$ for some integers $k,m$. Now, 
\begin{align*}
\dfrac{a}{b} - \dfrac{c}{d}= \dfrac{ad-bc}{bd}
\end{align*}

Since an odd integer times an odd integer is always odd, we get that $\dfrac{a}{b}-\dfrac{c}{d} \in \mathbb{Q'}$. Also, clearly, $1 \in \mathbb{Q'} \text{and } 0 \in \mathbb{Q'}$. So, $\mathbb{Q'}$ is a subring of $\mathbb{Q}$\\

\textbf{Q2}

\textbf{Q3: } Show the center of a ring is a subring containing identity. Show the center of a division
ring is a field.\\

\textbf{Soln: } Let $R$ be a ring and let $Z(R)$ be the center of $R$. We want to show that $Z(R)$ is a subring of $R$ containing identity. Let $a,b \in Z(R)$. Clearly, 1 and 0 are in $Z(R)$ since $1\cdot x=x\cdot 1$ for all $x \in Z(R)$, and similarly for 0 . Now let $c$ be an arbitrary element of $R$. Then, we have $(a-b)c=ac-bc=ca-cb=c(a-b)$. This means that $a-b \in Z(R)$ and thus $Z(R)$ is a subring of $R$ containing identity. \\

For the second part, we know that a division ring is a ring(not necessarily commutative) such that every element is invertible. If $R$ is a division ring, then $Z(R)$ is a subring of $R$ which is commutative, thus the result follows.

\textbf{Q4} If $R$ is an integral domain such that $x^2=1$, then $x=\pm1$.\\

\textbf{Soln: } $x^2=1\implies x^2-1=0\implies (x+1)(x-1)=0\implies x+1=0 \text{ or } x-1=0 \implies x=\pm 1$

\end{document}